\chapter{Transport, sélection et prolongement dans le noyau topologique de phase}
\label{sec:arquitectura-tpk}

Le sigle TPK désigne le \emph{Topological Phase Kernel}, ou noyau
topologique de phase.\footnote{La dénomination conserve également la
dédicace de l'auteur à Tan Peng Kian.} Il ne désigne pas un unique algorithme de 108
pas, un tableau de résultats ni le catalogue de 468 émissions. C'est le
système d'opérateurs qui fait coexister, sur la catégorie des chemins
d'APP, les lectures additive et multiplicative sans identifier leurs généalogies.

Le mot \emph{topologique} a ici un contenu précis. Le support
d'APP est le graphe de Cayley du groupe quotient
\[
 X_9=(\mathbb Z/9\mathbb Z)^2
\]
avec ses arêtes cardinales. Une voie fermée admet une classe d'enroulement
après relèvement au revêtement \(\mathbb Z^2\). TPK transporte cette classe,
l'orientation de la voie et l'holonomie de phase. Deux chemins qui se terminent
dans la même cellule peuvent donc représenter des états distincts. Le noyau
de phase est précisément l'information qui reste invisible lorsque l'on ne publie
que la cellule finale ou son reste.

La base dynamique est la catégorie libre \(\mathsf P_{\rm APP}\) du graphe
cardinal d'APP\@. Puisque les deux évaluations arithmétiques n'induisent pas la
même partition, le transport vit sur la catégorie bivariée
\[
 \mathsf P_{\rm APP}^{(2)}
 =\mathsf P_{\rm APP}\times\mathsf P_{\rm APP}.
\]
La première composante est évaluée au moyen du feuillet additif \(\Sigma\) et la
seconde au moyen du feuillet multiplicatif \(\Pi\). Cette séparation est
constitutive, mais ne duplique pas APP : elle enregistre deux lectures ordonnées d'une
même matrice de chemins. La composante multiplicative est le tableau principal ;
l'additive conserve la loi d'accumulation qui engendre ses lignes et permet de
mesurer l'information éliminée par une projection visible.

\ifdefined\HMTTPKRepeatAPP
\section{Carte canonique et architecture somme--produit}

Le support s'écrit en coordonnées de curseur
\[
 I_9=\{0,\ldots,8\},\qquad X_9=I_9^2,
\]
mais ses étiquettes arithmétiques appartiennent à \(\{1,\ldots,9\}\). La carte
qui relie les deux niveaux est
\[
 \chi:I_9\longrightarrow\{1,\ldots,9\},\qquad \chi(i)=i+1.
\]
Soit
\[
 \rho_9(n)=1+((n-1)\bmod9),\qquad
 q_9(n)=\frac{n-\rho_9(n)}9.
\]
En particulier, \(9\) est le représentant visible de la classe nulle. Les deux
observables d'un même sommet sont
\[
 \Sigma(i,j)=\rho_9\bigl(\chi(i)+\chi(j)\bigr),\qquad
 \Pi(i,j)=\rho_9\bigl(\chi(i)\chi(j)\bigr).
\]
Cette convention empêche de mélanger la coordonnée de curseur zéro avec l'étiquette
visible neuf. APP est le tuple
\[
 \mathcal A_9=
 \bigl(X_9,\operatorname{Cay}(X_9),\Sigma,\Pi,
       \operatorname{Path}(X_9)\bigr),
\]
et non la juxtaposition de deux tableaux : chaque germe \((i,j)\) porte à la fois
les deux évaluations et chaque voie conserve l'ordre dans lequel elle les consulte.

Le relèvement qui retient le reste et le quotient est
\[
 \mathcal A_9^\#=(R^+,Q^+;R^\times,Q^\times),
\]
avec, pour \(a=\chi(i)\) et \(b=\chi(j)\),
\[
 a+b=R^+_{ij}+9Q^+_{ij},\qquad
 ab=R^\times_{ij}+9Q^\times_{ij},
\]
\[
 R^+_{ij}=\Sigma(i,j),\qquad R^\times_{ij}=\Pi(i,j).
\]
Les totaux exacts sont
\[
 \sum R^+=405,\quad \sum R^\times=459,\quad
 \sum Q^+=45,\quad \sum Q^\times=174.
\]
Par conséquent,
\[
 2025-810=1215=(459-405)+9(174-45)=54+9\cdot129.
\]
L'égalité ne crée pas TPK par simple comptage : elle identifie la donnée que le
transport doit conserver. Le reste visible diffère de \(54\) et la mémoire
du quotient de \(129\).

La composition du feuillet additif possède le cocycle de retenue
\[
 \kappa_9(a,b)=
 \frac{\rho_9(a)+\rho_9(b)-\rho_9(a+b)}9,
\]
qui satisfait
\[
 \kappa_9(a,b)+\kappa_9(a+b,c)
 =\kappa_9(b,c)+\kappa_9(a,b+c).
\]
Le secteur multiplicatif possède en outre le radical
\[
 \mathfrak m=(3)=\{3,6,9\},\qquad \mathfrak m^2=0
 \quad\text{dans }\mathbb Z/9\mathbb Z.
\]
Les six lignes unitaires ont pour somme \(45\), les lignes \(3,6\) ont pour somme \(54\) et
la ligne \(9\) a pour somme \(81\) ; ainsi
\[
 459=6\cdot45+2\cdot54+81,
\]
et tout l'excès \(459-405=54\) se situe dans le secteur non unitaire. Cette
localisation, le cocycle de retenue et la différence des quotients sont trois
données distinctes de la même architecture somme--produit.
\fi

\section{Réductions résiduelles et représentation décimale}

L'architecture utilise trois réductions, chacune ayant un rôle distinct :
\[
 \mathbb Z\xrightarrow{(\rho_9,q_9)}
 \{1,\ldots,9\}\times\mathbb Z,
 \qquad
 \rho_3^{\rm or}:\mathbb Z/3\mathbb Z\longrightarrow
 \{-1,0,+1\},
\]
\[
 (d_1,d_2,\ldots)\in\{0,\ldots,999\}^{\mathbb N}
 \xrightarrow{\iota_{1000}}
 \sum_{m\geq1}d_m1000^{-m}.
\]
L'orientation ternaire est
\[
 \rho_3^{\rm or}([0])=0,\qquad
 \rho_3^{\rm or}([1])=+1,\qquad
 \rho_3^{\rm or}([2])=-1.
\]
La première sépare le reste visible et le quotient nonadique ; la deuxième attribue le
type orienté de la transition ; la troisième regroupe la lecture archimédienne
en triades décimales, dont le registre conserve la retenue. Aucune ne peut
remplacer une autre.
En particulier, la réduction modulo trois ne reconstruit pas le quotient modulo
neuf, et une triade décimale n'est pas un mot de six trits.
La somme accumule les déplacements ; le produit replie les classes et contracte les
lignes non inversibles. La nécessité de TPK naît du fait que les deux lectures
partagent le support, mais ne partagent ni la partition, ni la période, ni le quotient.

\Needspace{36\baselineskip}
\section{Correspondance entre symboles du TPK, domaines et compositions}
\label{sec:diccionario-interno-tpk}

Les noms historiques sont fixés ici par objet et non par ressemblance
numérique. Ce tableau fait partie de la définition publique de TPK.

\begin{longtable}{@{}>{\raggedright\arraybackslash}p{.19\textwidth}
>{\raggedright\arraybackslash}p{.30\textwidth}
>{\raggedright\arraybackslash}p{.43\textwidth}@{}}
\toprule
Nom&Objet&Fonction au sein du noyau\\\midrule
\endhead
coordonnée nonadique&
\(\rho_9(n)\) avec \(q_9(n)\)&
Sépare le représentant visible modulo \(9\) du quotient qui conserve la
mémoire du franchissement de frontière. Ce n'est pas une racine numérique.\\
sélecteur temporel ternaire&
\(M_{\rm ph}(t)=\varepsilon(t)\in\{+1,-1,0\}\)&
En \(1,4,7\), il active l'accumulation additive ; en \(2,5,8\), il active
l'accumulation multiplicative ; en \(3,6,9\), il maintient les deux curseurs et enregistre le seuil.
Il ne désigne pas les directions spatiales.\\
transporte cardinal&
\(T_N,T_E,T_S,T_O\)&
Déplace la position sur le tore APP et conserve le mot de voie et son
enroulement. Il est indépendant du sélecteur temporel.\\
lecteur ternaire&
\(\rho_3^{\rm or}\)&
Convertit les trois classes résiduelles en orientation
\(0,+1,-1\). Il ne reconstruit pas à lui seul le quotient nonadique.\\
lecteur en base mille&
\(\iota_{1000}\)&
Regroupe la publication archimédienne en triades décimales et transporte ses
retenues. C'est un lecteur positionnel, non une congruence ``modulo \(1000\)''.\\
retour de vingt-sept&
\(\mathcal K_{27}=F_{\rm obs}^{27}\)&
Compose vingt-sept mises à jour positionnelles et orientées. Son quatrième
itéré restaure la phase observable, tandis que le registre complet conserve
la mémoire accumulée.\\
registre dodécaphasique&
\(\Lambda\in\mathcal R_{12}\)&
Conserve six positions, cinq différences et un total. C'est le support du
transport réversible \(1+5+6=12\).\\
lecture bidirectionnelle&
\(q_+,q_-\) et ses poids conjugués&
Conserve simultanément les orientations de l'aller et du retour. En mécanique
dimensionnelle, une signature \((n_A,s_C)\) produit d'abord
\(W_\pm=6n_A\pm s_C\) ; ce n'est qu'ensuite qu'apparaît
\(s_C/6\) dans le caractère angulaire.\\
\bottomrule
\end{longtable}

Le sélecteur temporel, le transport cardinal et l'orientation ternaire sont
donc trois applications distinctes. Leur séparation conserve la différence entre
\emph{quand} une lecture est activée, \emph{vers où} la voie avance et
\emph{avec quelle orientation} le pas est enregistré. TPK contient les trois et
conserve leurs relations.

\ifdefined\HMTTPKWordArithmetic
\section{De la cellule locale au monoïde des mots nonadiques}

L'architecture somme--produit ne s'arrête pas au tableau local. Pour prouver que
le quotient transporté engendre une arithmétique globale, il faut construire
l'opération sur les mots sans introduire comme définition la multiplication
entière que l'on entend retrouver. Soit
\[
 \mathcal D_9=\{1,\ldots,9\},\qquad
 \mathcal W_9=\{\varnothing\}\sqcup
 \bigsqcup_{\ell\geq1}\mathcal D_9^\ell .
\]
Les mots s'écrivent à partir de la position la moins significative et s'évaluent
par
\[
 \operatorname{ev}_9(u_0,\ldots,u_{\ell-1})
 =\sum_{j=0}^{\ell-1}u_j9^j,\qquad
 \operatorname{ev}_9(\varnothing)=0.
\]
Le symbole \(9\) conserve le franchissement de frontière ; il n'est pas remplacé par un zéro.

\begin{theorem}[Forme normale nonadique]
L'évaluation \(\operatorname{ev}_9\) est une bijection entre
\(\mathcal W_9\) et \(\mathbb N_0\).
\end{theorem}
\begin{proof}
Les mots de longueur \(\ell\) couvrent exactement l'intervalle
\[
 \frac{9^\ell-1}{8}
 \leq\operatorname{ev}_9(u)
 \leq9\frac{9^\ell-1}{8}.
\]
Le minimum de l'intervalle suivant dépasse d'une unité le maximum du
précédent. Pour \(n>0\), la division de \(n-1\) par neuf détermine de manière
unique
\[
 d=1+((n-1)\bmod9)\in\mathcal D_9
\]
et un quotient strictement inférieur, auquel la même règle s'applique
récursivement. Le processus se termine et est réversible.
\end{proof}
Par exemple,
\[
 9\leftrightarrow(9),\quad
 10\leftrightarrow(1,1),\quad
 81\leftrightarrow(9,8),\quad
 1000\leftrightarrow(1,3,3,1).
\]
La dernière identité exhibe dans une seule forme normale les deux partitions
fondamentales
\[
 1000=729+271=729+243+27+1;
\]
leur interprétation par canaux est introduite ensuite, non dans la définition
de la numération.

\section{Somme, produit de Horner et conservation du quotient}

La somme native \(u\boxplus_\#v\) parcourt les mots par positions. À chaque
position, elle applique la cellule additive aux deux chiffres présents et, s'il existe
une retenue entrante, applique une seconde cellule au reste provisoire. Le
reste est publié et la somme des quotients passe à la position suivante.
L'invariant partiel est
\[
 \sum_{j<k}(u_j+v_j)9^j
 =\sum_{j<k}w_j9^j+c_k9^k,
\]
d'où l'on obtient
\[
 \operatorname{ev}_9(u\boxplus_\#v)
 =\operatorname{ev}_9(u)+\operatorname{ev}_9(v).
\]

Pour \(d\in\mathcal D_9\), le transducteur scalaire
\(d\odot_\#u\) utilise d'abord
\[
 u_jd=9q_j+r_j,\qquad
 (r_j,q_j)=\bigl(R^\times(u_j,d),Q^\times(u_j,d)\bigr),
\]
et normalise la collision entre \(r_j\) et le quotient entrant au moyen de la
cellule additive. Ainsi,
\[
 \operatorname{ev}_9(d\odot_\#u)
 =d\operatorname{ev}_9(u).
\]
Si \(v=(v_0,\ldots,v_{m-1})\), on définit, de droite à gauche,
\[
 a_m=\varnothing,\qquad
 a_j=(9\odot_\#a_{j+1})\boxplus_\#(v_j\odot_\#u),
\]
et
\[
 u\boxtimes_\#v:=a_0.
\]
C'est une construction de Horner réalisée uniquement avec les deux cellules
locales et leurs quotients.

Cette description fixe la dépendance opératoire, mais ne constitue pas un
second emplacement du résultat. Le produit global, sa preuve de Horner et
l'identité
\(operatorname{ev}_9(u\boxtimes_\#v)
=\operatorname{ev}_9(u)\operatorname{ev}_9(v)\)
sont établis une seule fois dans le
\cref{int:eq:ev-multiplicative}, au sein du chapitre spécialisé sur le produit
natif.
L'exemple de frontière le plus court est
\[
 (9)\boxtimes_\#(9)=(9,8),\qquad 9+8\cdot9=81.
\]
Ne conserver que \(R^\times\) détruirait ce théorème : \(2\cdot5=10\)
serait réduit au reste visible \(1\). Le quotient n'est pas une annotation
auxiliaire, mais la mémoire nécessaire pour que le produit survive au passage de
la frontière.

\section{Coproduit et irréductibilité}

L'ouverture multiplicative est définie de manière interne par
\[
 \Delta_\#e_w
 =\sum_{u\boxtimes_\#v=w}e_u\otimes e_v.
\]
Sa version réduite élimine les facteurs \((1)\). Si
\(\widetilde\Delta_\#\) désigne cette version et
\[
 A_\#=
 \bigl(\overline{\widetilde\Delta_\#}\bigr)^*
 \overline{\widetilde\Delta_\#},
\]
alors \(A_\#e_w=d_\#(w)e_w\), où \(d_\#(w)\) compte les paires ordonnées
non triviales qui produisent \(w\). Le sélecteur
\[
 P_{\mathrm{Irr},\#}
 =(I-P_{(1)})\,\mathbf1_{\{0\}}(A_\#)
\]
choisit exactement les mots distincts de l'unité sans ouverture native non
triviale. L'irréductibilité se lit donc dans la fibre même du produit ;
elle n'est pas introduite au moyen d'une table externe de nombres premiers.

Ce résultat est exact dans le monoïde des formes normales. Son extension aux
histoires complètes exige une compatibilité supplémentaire. Si \(P_n\) projette
sur des histoires survivantes de profondeur \(n\), une famille
d'injections \(\iota_n\) doit satisfaire
\[
 P_n\iota_n\mu_\#
 =P_n\mu_n(\iota_n\otimes\iota_n).
\]
Les théorèmes précédents construisent \(\mu_\#\) et fixent cette équation ; ils
n'affirment pas à eux seuls l'existence ou l'unicité de la famille
\((\iota_n)_n\). Cette séparation empêche de confondre l'arithmétique globale des
mots, déjà démontrée, avec son prolongement monoïdal à l'espace profond de
mémoire TPK\@.

\fi

\section{Action cardinale et rétroaction arithmétique}

Soient
\[
 \nu_N=(-1,0),\quad \nu_E=(0,1),\quad
 \nu_S=(1,0),\quad \nu_O=(0,-1)
\]
les quatre générateurs cardinaux. Pour \(d\in\{N,E,S,O\}\),
\[
 T_d(i,j)=(i,j)+\nu_d\pmod9.
\]
La lettre \(O\) signifie ouest. Ces directions sont des déplacements sur
APP et non des noms des canaux \(\pi,e,\varphi\).

Un mot \(w=d_1\cdots d_r\) détermine
\[
 \gamma(x_0;w)_k=x_0+\sum_{j=1}^{k}\nu_{d_j}\pmod9.
\]
Lorsqu'il se ferme, son relèvement conserve
\[
 \operatorname{wind}(w)=\frac1{9}
 \bigl(\#S-\#N,\#E-\#O\bigr)\in\mathbb Z^2.
\]
La lecture bidirectionnelle conserve le mot, l'orientation et le feuillet. Son
renversement formel inverse l'ordre et échange
\(N\leftrightarrow S\), \(E\leftrightarrow O\). Son éventuelle lecture comme
avance--retard est une interface physique ultérieure, non une partie de cette
définition de voie.

L'état local observable est
\[
 s_t=(i_t,j_t,m_t),
 \qquad m_t\in\{\Sigma,\Pi\}.
\]
On évalue le feuillet actif,
\[
 a_t=\rho_3^{\rm or}\bigl(m_t(i_t,j_t)\bigr),
\]
et le mode est mis à jour au moyen de
\[
 m_{t+1}=
 \begin{cases}
 \Sigma,&a_t=+1,\\
 \Pi,&a_t=-1,\\
 m_t,&a_t=0.
 \end{cases}
\]
Ainsi, chaque pas cardinal modifie à la fois la position et la lecture
arithmétique future. La voie ne se superpose pas à un mot déjà construit : c'est
l'histoire positionnelle qui le produit.

\begin{theorem}[Surjectivité du facteur échantillonné de pas trois]
Pour chacun des \(162=9^2\cdot2\) états \((i,j,m)\) et chaque trit
orienté cible \(b\in\{-1,0,+1\}\), il existe une voie
\(u\in\{N,E,S,O\}^3\) dont le troisième pas émet \(b\). Dans les \(486\) cas
état--cible, le nombre de voies admissibles est compris entre
\(8\) et \(40\). Par conséquent, pour toute suite cible
\((b_k)_{k\geq1}\), il existe une histoire cardinale dont la sortie satisfait
\[
 (a_3,a_6,a_9,\ldots)=(b_1,b_2,b_3,\ldots).
\]
Les deux trits intermédiaires \(a_{3k-2}\) et \(a_{3k-1}\) de chaque bloc ne
sont pas prescrits par cet énoncé.
\end{theorem}
\begin{proof}
On énumère les \(4^3=64\) voies de trois pas à partir de chaque état. Le
certificat exhaustif inclus enregistre, pour chacun des
\(162\cdot3\) cas, la voie exécutée, l'état final et le nombre total de
témoins ; leurs indices \(0,1,2\) sont transportés par
\(\rho_3^{\rm or}\) vers \(0,+1,-1\). Tous les décomptes sont positifs ; le
minimum est \(8\) et le maximum est \(40\). Après chaque bloc, le même
résultat s'applique à l'état atteint.
Le choix, par exemple, du premier témoin dans l'ordre lexicographique
produit récursivement une histoire pour toute suite cible.
\end{proof}

La surjectivité concerne le facteur observé tous les trois pas, non le
mot émis au temps unitaire. Ce théorème établit cette capacité
échantillonnée du support des voies. La sélection des sections
\(\pi,e,\varphi\) appartient à la relation enrichie EF--G9 définie
plus loin : la capacité de réalisation et la sélection généalogique sont deux
opérateurs distincts.

\section{Chronologie observable des deux curseurs}

Le calendrier se formule sur
\[
 \Theta_{108}=\mathbb Z/108\mathbb Z,\qquad
 \beta(\theta)=
 \left[\left\lfloor\frac{\widetilde\theta}{9}\right\rfloor\right]_{12},
 \qquad
 r(\theta)=1+(\theta\bmod9),
\]
où \(\widetilde\theta\in\{0,\ldots,107\}\). La signature active est
\[
 \varepsilon(r)=
 \begin{cases}
 +1,&r\in\{1,4,7\},\\
 -1,&r\in\{2,5,8\},\\
 0,&r\in\{3,6,9\}.
 \end{cases}
\]
L'espace observable à deux curseurs est
\[
 \mathcal Q_{\rm obs}
 =X_9^2\times\{N,E,S,O\}^2\times\Theta_{108}.
\]
Si \(\varepsilon=+1\), le curseur additif avance ; si
\(\varepsilon=-1\), le multiplicatif ; si \(\varepsilon=0\), aucun.
Les deux orientations s'inversent après les pas
\(27,54,81,108\). Cette règle définit une permutation
\[
 F_{\rm obs}:\mathcal Q_{\rm obs}\longrightarrow\mathcal Q_{\rm obs}.
\]
Les positions et orientations sont restaurées après \(54\) pas et la phase
complète après \(108\) :
\[
 F_{\rm obs}^{108}=I.
\]
La chronologie d'une seule voie est donc réversible et ne sélectionne pas à elle
seule une distribution sur les continuations ; la sélection de tous les
prolongements compatibles appartient à l'opérateur de transfert des fibres
défini ci-dessous.

\section{Émetteur fini à deux curseurs}

Un germe de curseur est
\[
 \mathcal S=\mathbb T_9^2\times\{N,E,S,O\},
 \qquad |\mathcal S|=324.
\]
TPK minimal reçoit un curseur additif \(P\in\mathcal S\) et un curseur
multiplicatif \(M\in\mathcal S\). Leurs signatures de six fenêtres sont
\(s(P)\) et \(e(M)\). À chaque coordonnée, on applique
\[
 G(s,e)=100s+10\bigl((s+e)\bmod10\bigr)
       +\bigl((3s+5e+1)\bmod10\bigr).
\]
Le chiffre des centaines retrouve \(s\) ; \(s\) étant connue, celui des dizaines retrouve
\(e\). Ainsi \(G\) est injective sur les signatures effectives. L'émetteur
se factorise comme
\[
 \mathcal S^2\xrightarrow{s\times e}
 \operatorname{im}s\times\operatorname{im}e
 \xrightarrow{G}\mathcal U_6,
\]
avec
\[
 |\mathcal S^2|=104\,976,\qquad
 |\operatorname{im}s|=18,\qquad
 |\operatorname{im}e|=26,\qquad
 |\mathcal U_6|=18\cdot26=468.
\]
La factorisation \(468=18\cdot26\) est donc une décomposition de
l'objet : chaque émission conserve une signature additive et une multiplicative.

La projection
\[
 \Psi_6:\mathcal U_6\longrightarrow\mathbb F_3^6,
 \qquad \Psi_6(U)=-U\pmod3,
\]
possède une image de \(243\) mots. La chaîne typée est
\[
 104\,976\ \text{germes}
 \longrightarrow468\ \text{émissions décimales }U_6
 \longrightarrow243\ \text{mots ternaires }w_6.
\]
L'espace ambiant \(\mathbb F_3^6\) contient \(729\) mots ; il ne s'identifie pas
à l'image effective de \(243\). \(U_6\) et \(w_6\) ne sont pas non plus de même
type.

L'émetteur minimal répète, dans la carte dodécaphasique du catalogue,
\[
 U_{12}^{\rm cat}=U_6\Vert U_6.
\]
Sa période visible six explique pourquoi il ne peut pas remplacer l'état signé
\(U_{12}^{\rm sgn}\). Cette différence n'est pas une insuffisance de TPK : elle sépare
la projection de catalogue de l'histoire enrichie.

Chaque bloc décimal émis appartient à un alphabet exact de
\(7\cdot6=42\) triades. Si \(\overline{d_1d_2d_3}\) est l'une d'elles,
\[
 \boxed{d_3\equiv5d_2-2d_1+1\pmod{10}.}
\]
Le balayage des \(104\,976\) germes ne produit aucune triade en dehors de cette
loi. Il convient donc de distinguer la chaîne complète des quotients :
\[
 \mathcal S^2
 \longrightarrow\Omega_{\rm seed}
 \longrightarrow\mathcal U_6
 \xrightarrow{\Psi_6}\operatorname{im}\Psi_6\subset\mathbb F_3^6.
\]
Les présentations microscopiques, les \(468\) émissions décimales, les
\(243\) mots ternaires et l'espace ambiant de \(729\) mots sont quatre
objets distincts.

\paragraph{Commutativité orbitale des morphismes typés.}
Soit
\(W_{243}:=\operatorname{im}\Psi_6\subset\mathbb F_3^6\). L'entonnoir
complet et son quotient diédrique se représentent sans transformer les inclusions en
sélections :
\[
 \underbrace{\mathcal S^2}_{104\,976\ \mathrm{semillas}}
 \xrightarrow[\text{recensement des fibres}]{\text{émission}}
 \underbrace{\mathcal U_6}_{468\ U_6}
 \xrightarrow[\text{projection visible}]{\Psi_6}
 \underbrace{W_{243}}_{243\ w_6}
 \hookrightarrow
 \underbrace{\mathbb F_3^6}_{729\ \mathrm{palabras}},
 \qquad
 W_{243}\xrightarrow{\,/D_3^{\rm cat}\,}
 \underbrace{\mathcal O_{43}}_{43\ \text{orbites}}.
 \tag{TPK-embudo}\label{eq:embudo-orbital-tipado-rev8}
\]
La première flèche regroupe les conditions initiales selon l'émission qu'elles produisent ;
la deuxième oublie la présentation décimale ; la troisième est une inclusion dans
l'espace ambiant algébrique ; la dernière est un quotient par l'action
\(D_3^{\rm cat}\) sur \(W_{243}\). En particulier, il n'y a pas de quotient
\(729\to43\) implicite. Le recensement exact est
\[
 243=38\cdot6+5\cdot3,
\]
c'est-à-dire trente-huit orbites de taille six et cinq de taille trois.
Ces \(43\) classes orbitales ne sont pas non plus les \(43\) tours nonadiques
complets par canal du témoin G9 de mille chiffres : le premier \(43\) appartient
au quotient fini \(W_{243}/D_3^{\rm cat}\) ; le second est un compteur de
profondeur d'une exécution certifiée.

\section{Configuration de l'état enrichi du TPK : feuillet, orientation, reste, quotient, retenue, frontière, voie et mémoire}

L'unité dynamique complète est une histoire, non un mot isolé. À la
profondeur \(m\), nous écrivons
\[
 v_m=
 \bigl(
 B_m,\boldsymbol\delta_m,\boldsymbol x_m,
 \boldsymbol{\mathcal C}_m,\Sigma_m,
 \omega_m,\ell_m,g(m),\Lambda_m
 \bigr),
\]
où
\[
 B_m=(u_m^\pi,u_m^e,u_m^\varphi)\in(\mathbb F_3^6)^3.
\]
Ses composantes ont des fonctions distinctes :
\begin{enumerate}
 \item \(B_m\) est le bloc visible conjoint des trois canaux ;
 \item \(\boldsymbol\delta_m\) conserve les nouvelles triades décimales et la
       retenue associée ;
 \item \(\boldsymbol x_m\) est l'état de Hensel profond ;
 \item \(\boldsymbol{\mathcal C}_m\) contient les cylindres ternaires et
       décimaux et leurs restes entiers ;
 \item \(\Sigma_m\) est la signature conjointe de frontière ;
 \item \(\omega_m\) et \(\ell_m\) conservent l'orientation et le feuillet bidirectionnel ;
 \item \(g(m)\in\mathbb Z/9\mathbb Z\) est la phase nonadique ;
 \item \(\Lambda_m\) est le registre dodécaphasique orienté.
\end{enumerate}

\begin{definition}[Présentation opératoire et catégorie TPK]
L'inventaire des opérateurs de l'état enrichi de TPK est
\[
 \mathfrak T_{\rm HMT}=
 \bigl(
 \mathcal A_9,\TRIT,T_{N,E,S,O},T_{\min},
 H,\mathcal C_{3,1000},\Sigma_B,EF,\mathcal G_9,
 \mathcal R_{12}
 \bigr)
\]
qui transporte les histoires \((v_0,\ldots,v_m)\) et leurs troncatures. Son
support formel est la catégorie \(\mathsf{TPK}\) sur la base observable :
ses objets sont des états enrichis et ses morphismes sont les prolongements
\(\operatorname{Ext}_r\) stables par identité et concaténation, définis
ci-dessous. L'inventaire ne remplace pas cette catégorie. La projection
visible élimine des composantes de \(v_m\) ; elle ne définit pas d'équivalence avec
l'état complet.
\end{definition}

Cette définition contient exactement ce qui disparaît lorsque les curseurs sont réinitialisés
ou que seul le reste est conservé : le chemin, le feuillet, le quotient, la
retenue, l'orientation et la frontière. TPK est donc le point où
l'arithmétique modulaire cesse d'être une opération sans histoire.

\section{Cylindres ternaires-décimaux et état d'extension}

Soit \(N\) l'entier déterminé par un préfixe ternaire de longueur \(L\), et
soit \(D\) l'entier déterminé par \(K\) triades décimales. Nous définissons
\[
 A=N1000^K-D3^L,
 \qquad
 B=(N+1)1000^K-D3^L.
\]
Les cylindres se coupent exactement lorsque
\[
 \boxed{A<3^L,\qquad B>0.}
\]
En incorporant un bloc ternaire \(u\in\{0,\ldots,728\}\),
\[
 N'=729N+u.
\]
S'il n'apparaît pas de nouvelle triade décimale,
\[
 A'=729A+u1000^K.
\]
Si aparece \(\delta\in\{0,\ldots,999\}\),
\[
 D'=1000D+\delta,
\]
\[
 A'=729000A+u1000^{K+1}-\delta\,729\,3^L.
\]
Dans les deux cas,
\[
 B'=A'+1000^{K+\Delta K}.
\]
Ces récurrences entières transportent la donnée décimale sans en faire un
objectif extérieur de la mise à jour.

L'état de Hensel conserve l'information qui n'apparaît pas dans le reste. La
matrice active est
\[
A_H=\begin{pmatrix}
2&2&2&1&2&1\\2&1&2&2&1&1\\1&0&0&1&0&2\\
0&0&1&1&1&0\\2&2&2&0&2&1\\2&1&2&1&0&0
\end{pmatrix},\qquad \det A_H=1,
\]
et la division reste--quotient est
\[
 y=xA_H,\qquad u=y\bmod3,\qquad x^+=\frac{y-u}{3}.
\]
Pour les trois canaux, on forme en outre les cocycles conjoints
\[
 q=(\pi+e+\varphi)\bmod3,
 \qquad
 c=\frac{\pi+e+\varphi-q}{3},
\]
\[
 a=(\pi+e-\varphi)\bmod3,
 \qquad
 h=\frac{\pi+e-\varphi-a}{3}.
\]
La mise à jour du second ordre corrige l'état suivant avec \(c\) et
\(h\). C'est pourquoi deux branches ayant le même bloc visible ne sont pas nécessairement le
même état : elles peuvent différer par leur extension de Hensel et par leurs cylindres.
Comme \(\det A_H=1\), la réduction de \(A_H\) appartient à
\(\operatorname{GL}_6(\mathbb F_3)\). Ainsi, avant d'imposer les cylindres,
les frontières ou la survie, la carte résiduelle de Hensel couvre exactement
\(\mathbb F_3^6\). Cette couverture brute n'implique pas la couverture du système
filtré.

Formellement, sur chaque configuration observable \(q\), on fixe une fibre
\[
 \mathcal F_q=
 \mathsf O_q\times\mathsf H_q\times\mathsf C_q\times
 \mathsf S_q\times\mathsf B_q\times\mathsf\Lambda_q,
\]
correspondant, respectivement, à l'orientation et au feuillet, au reste et au quotient,
à la retenue, à la signature, au cylindre et à la frontière, et au registre de transport. Pour une
flèche cardinale \(r:q\to q'\), le prolongement admissible ne se réduit pas à une
fonction sur les restes ; c'est la relation
\[
 \operatorname{Ext}_r\subseteq\mathcal F_q\times\mathcal F_{q'},
 \qquad
 \operatorname{Ext}_{r_2}\circ\operatorname{Ext}_{r_1}
 \subseteq\operatorname{Ext}_{r_2\circ r_1}.
\]
L'incrément de retenue est un cocycle :
\[
 \kappa(r_2\circ r_1)
 =\kappa(r_2)+\mathsf C(r_2)\bigl(\kappa(r_1)\bigr).
\]
L'application d'oubli est ainsi typée : projeter une extension sur sa voie
observable élimine des composantes de \(\mathcal F_q\), mais n'autorise pas à
identifier les deux objets.

\section{Retours composés et calendriers de transport}

La longueur d'une trace ne détermine pas son type. La présentation académique
sépare le retour composé, le calendrier, la trace et l'atlas selon leurs
états conservés respectifs.

Le symbole \(\mathcal K_{27}\) est réservé à la composition de
vingt-sept mises à jour d'un lacet de transport positionnel et orienté.
L'entier \(27\) est sa longueur de composition ; il n'en fait pas une
trace de 108 itérations ni un émetteur infini. Une autre construction historique
regroupe quatre lacets dont les périodes conjointes atteignent 108 ; ce rotor conjoint
n'est pas désigné par \(\mathcal K_{27}\) isolé.

\begin{longtable}{@{}>{\raggedright\arraybackslash}p{.24\textwidth}>{\raggedright\arraybackslash}p{.31\textwidth}>{\raggedright\arraybackslash}p{.37\textwidth}@{}}
\toprule
Objet&Estado conservado&Significado de 108\\\midrule
\endhead
Contador simplificado&position et vitesse&sa période réelle est 18 ; c'est un
contrôle historique.\\
Rotor conjoint de quatre lacets&quatre voies, lecture, rotation et mouvement&
les périodes individuelles sont \(54,54,36,36\) ; l'état conjoint a une
période de 108.\\
Trace brute cadencée&quatre positions, orientations et phase&le contenu
dynamique se répète à 54.\\
Trace de \(T_{\min}\)&deux positions et deux directions&l'état revient à
54 et la sortie est \(U_6\Vert U_6\).\\
Pseudocode réinitialisé&curseurs réinitialisés à chaque fenêtre&s'effondre en
\((222,222,222,333,333,333)^2\) ; ne représente pas l'état enrichi de TPK.\\
Rutas condicionadas&chemin et mot prescrit&108 est un horizon, non une période.\\
Odomètre sectoriel&\(\mathbb Z/9\mathbb Z\times\{0,\ldots,11\}\)&l'état
complet a une période de 108 ; l'observable sectoriel, de 36.\\
Projection brute&quotient par oubli de \(X_{\rm TPK}^{\rm enr}\)&ne possède pas
à elle seule la lecture signée.\\
Traza enriquecida&tous les champs de \(v_m\) pendant 108 mises à jour&est
la restriction temporelle de l'état enrichi de TPK.\\
Atlas des voies&81 cellules/56 signatures ou 324 germes orientés&c'est une taxonomie,
non une orbite de 108 états.\\
\bottomrule
\end{longtable}

Le lacet \(\mathcal K_{27}\) est l'opérateur de transport fini de
vingt-sept mises à jour. Les fenêtres, le rotor conjoint et la projection
calendaire reçoivent des symboles distincts, de sorte que la longueur partagée
ne remplace ni le domaine ni la composition de chaque objet.

\section{Transfert des fibres et mémoire de rotation}

La factorisation \(\mathcal U_6\simeq\mathcal S_+\times\mathcal S_\times\)
permet de définir, par rapport au comptage uniforme, les espérances conditionnelles
\[
 (E_+f)(s,e)=\frac1{18}\sum_{s'\in\mathcal S_+}f(s',e),
 \qquad
 (E_\times f)(s,e)=\frac1{26}
 \sum_{e'\in\mathcal S_\times}f(s,e').
\]
\(E_+\) renouvelle la fibre additive et conserve la multiplicative ;
\(E_\times\) réalise l'opération duale. Ce sont les seuls projecteurs de Markov
qui marginalisent complètement la coordonnée active, conservent l'inactive et
préservent le comptage uniforme.

Soit \(\tau=(+1,-1,0)^3\) le mot de phase sur \(C_9\), et
\(P_{+1}=E_+\), \(P_{-1}=E_\times\), \(P_0=I\). Le transfert nonadique est
\[
 \mathcal K_{\rm ph}\bigl((u,r),(v,r+1)\bigr)
 =P_{\tau(r)}(u,v).
\]
Il est doublement stochastique, irréductible, de période \(9\), et possède une unique
distribution stationnaire \(1/(9\cdot468)\). De plus,
\[
 \mathcal K_{\rm ph}^{\,9}=E_+E_\times,\qquad
 (E_+E_\times)(u,v)=\frac1{468}.
\]
Cette égalité exprime le renouvellement simultané de toutes les continuations
compatibles, non le successeur déterministe d'une seule voie.

Trois opérations géométriques sur les germes conservent une information
supplémentaire : \(P\), avancée dans l'orientation active ; \(H\), demi-tour ; et
\(Q\), quart de tour. La paire \((P,H)\) raffine les \(26\) signatures produit
en \(28\) classes. En particulier, la signature \(555555\) se sépare en trois feuillets
de huit représentants, et
\[
 18\cdot28=504=468+36.
\]
Le terme \(36\) compte les feuillets supplémentaires du raffinement et n'est pas
l'objet \(R_{36}\). La paire \((P,Q)\) produit \(162\) classes de deux
éléments, permutées transitivement par
\[
 J(i,j,d)=\bigl(7-i,\,7-j,\,\overline d\bigr)\pmod9.
\]

Dans une trace de \(108\) pas apparaissent \(36\) signes de chaque type,
\(72\) changements effectifs entre les deux feuillets et quatre inversions
d'orientation ; le cocycle signé total est nul. Ces décomptes sont des invariants
du registre d'évolution et relient la dynamique nonadique au registre dodécaphasique.

\section{Prolongement ternaire et frontière conjointe}

Relativement aux germes, frontières et lecteurs nommés du profil
publié, les mots visibles initiaux sont
\[
 w_6^\pi=010211,\qquad w_6^e=201101,\qquad
 w_6^\varphi=121200.
\]
Après cinq blocs, les préfixes sont
\[
\begin{aligned}
w_{30}^{\pi}&=010211|012222|010211|002111|110221,\\
w_{30}^{e}&=201101|121221|102011|012222|102011,\\
w_{30}^{\varphi}&=121200|112202|121020|010210|010200.
\end{aligned}
\]
La première frontière profonde est
\[
 R_{36}=
 \begin{pmatrix}222220\\021222\\102011\end{pmatrix}.
\]
Pour chaque canal, on définit le préfixe
\(w_{36}^\chi:=w_{30}^\chi|u_5^\chi\) ; \(R_{36}\) est le triplet ordonné des
trois sixièmes blocs. Aucun n'est un émetteur ni un état terminal.

Pour \(B_k=(u_k^\pi,u_k^e,u_k^\varphi)\), les signatures
\[
 q_k=u_k^\pi+u_k^e+u_k^\varphi,\qquad
 a_k=u_k^\pi+u_k^e-u_k^\varphi
\]
déterminent
\[
 u_k^\varphi=2(q_k-a_k),\qquad
 u_k^\pi+u_k^e=q_k-u_k^\varphi
 \quad\text{dans }\mathbb F_3^6.
\]
Ainsi, \(\varphi\) est l'axe fixé et la liberté résiduelle est la répartition
interne \(\pi/e\). Dans la neuvième phase,
\[
 B_9=(100100|020112|010122)
\]
admet trois sorties locales ayant le même axe \(\varphi\) et le même agrégat
\(\pi+e\) ; les trois survivent à un horizon et seule
\[
 B_{10}=(101222|112221|112002)
\]
survit au second. C'est la descente exécutable \(3\to1\) d'EF--G9.
La branche publiée se poursuit jusqu'à vingt blocs par canal : \(w_{30}\) et
\(R_{36}\) sont préfixe et frontière, non clôture.

\section{Relation d'extension et monodromie nonadique}
\label{sec:extension-monodromia-tpk}

Soit \(\mathcal H_m\) l'ensemble des histoires enrichies admissibles de
profondeur \(m\), et soit
\[
 p_{n,m}:\mathcal H_n\longrightarrow\mathcal H_m
\]
la troncature. Pour \(N\geq m\),
\[
 \operatorname{Surv}_m^{(N)}=p_{N,m}(\mathcal H_N),
 \qquad
 \operatorname{Surv}_m=\bigcap_{N\geq m}\operatorname{Surv}_m^{(N)}.
\]
La relation EF--G9 conserve les extensions qui appartiennent à ce noyau
stable. La phase satisfait
\[
 g(m+9)=g(m),
\]
mais l'état transporté satisfait, en général,
\[
 v_{m+9}\ne v_m.
\]
Le retour agit sur la fibre des états d'une même phase ; il ne réinitialise pas
l'histoire.

\begin{definition}[Nombre HMT]
\label{def:numero-hmt-tpk-completo}
Un nombre HMT du canal \(\chi\) est uniquement une section globale
compatible
\[
 x=(x_m)_{m\geq0}\in
 X_\chi:=\varprojlim_m\operatorname{Surv}_m^{(\chi)}.
\]
Le lecteur et la valeur publiée ne font pas partie de la section.
\end{definition}

Une présentation ou mesure du nombre est une paire \((x,L)\), où
\(L\) est un lecteur déclaré dont le domaine contient \(x\). Lorsque \(L\) est le
lecteur archimédien des cylindres, \(L(x)\) est la valeur publiée ; \(x\)
conserve en outre l'orientation, le quotient, la retenue et la frontière.

\begin{theorem}[Existence d'un prolongement global]
\label{thm:prolongacion-global-tpk-completo}
Soit \(h\in\mathcal H_m\). Si l'arbre des prolongements de \(h\) est
à branchement fini et possède au moins un nœud à toute profondeur, alors
il existe une histoire infinie dont les troncatures prolongent \(h\). Si, de plus,
chaque niveau stable est un singleton et que les états uniques sont compatibles, la
section globale est unique.
\end{theorem}
\begin{proof}
Les niveaux finis non vides forment un arbre à branchement fini. Le
lemme de König fournit une branche infinie. La compatibilité des
troncatures identifie cette branche à un élément de la limite inverse. Si
chaque niveau contient un seul état, il n'existe pas de seconde branche compatible.
\end{proof}

Le théorème précédent est abstrait et conditionnel : il n'affirme pas à lui seul la
non-vacuité ni l'unicité stable d'un canal concret. Pour
\(\pi,e,\varphi\), ces hypothèses sont établies ensuite au moyen de la filtration
rationnelle exécutable du
\cref{thm:generacion-monodromica-conjunta-rev7}. Cet emplacement ne conserve
ici que le principe de limite inverse et n'anticipe pas son application.

\section{Composition des opérateurs fondamentaux}

Les opérateurs précédents conduisent à l'état terminal enrichi sans
aplatir ses coordonnées :
\[
 \begin{aligned}
 \mathrm{APP}&\longrightarrow\mathrm{TRIT}
 \longrightarrow\text{route orientée}
 \longrightarrow\text{état TPK enrichi}\\
 &\longrightarrow\text{cylindres et survie}
 \longrightarrow x_{\rm term}.
 \end{aligned}
\]
La fermeture dodécaphasique reçoit alors deux évaluations coordonnées du même
état :
\[
 x_{\rm term}
 \xrightarrow{(\pi_{\rm term},R_{\rm sgn})}
 \bigl(B_{\rm term},U_{12}^{\rm sgn}\bigr),
 \qquad
 U_{12}^{\rm sgn}\xleftrightarrow{\mathcal H_{12}}K,
\]
\[
 (P,E,\Phi,K,c_{13}=0)\longleftrightarrow\alpha_{12}.
\]
Il n'apparaît pas de flèche \(B_{\rm term}\to U_{12}^{\rm sgn}\) : les deux
coordonnées proviennent de \(x_{\rm term}\). Le passage ultérieur à la mécanique
dimensionnelle ne crée pas non plus une autre machine ; il prolonge la généalogie conservée
au moyen du registre chronologique enrichi, de la signature, du caractère et des tours.

\ifdefined\HMTInternalTraceability
\section{Registre académique et traçabilité}

Le fichier \path{datos/registro_operadores_tpk.json} fait partie du paquet.
Pour chaque opérateur, il déclare le nom académique, l'alias historique, le domaine,
le codomaine, la transition, les invariants, le statut et la source locale. Il inclut également
la séparation des dix constructions historiques associées au nombre
108. La porte de publication vérifie ce registre et bloque une sortie
qui réduirait de nouveau TPK à son catalogue ou à une trace brute.

Dans le reste de l'article, quatre abréviations, déjà typées, seront utilisées :
\[
 w_6\in\mathbb F_3^6,\qquad
 w_{30}\in\mathbb F_3^{30},\qquad
 R_{36}\in(\mathbb F_3^6)^3,\qquad
 U_{12}^{\rm sgn}\in\mathbb Z^{12}.
\]
La première est un mot germe ; la deuxième, un préfixe de cinq blocs ;
 la troisième, la frontière conjointe d'entrée dans EF--G9 ; la quatrième, la
coordonnée signée déclarée de l'état dodécaphasique enrichi. Aucune égalité de cardinalités n'autorise à
les identifier.
\fi
